\documentclass[12pt]{article}
\usepackage[utf8]{inputenc}
\usepackage[T1]{fontenc}
\usepackage{hyperref}
\title{Truk\'a Language Resources for K-12 Curriculum Development}
\author{Compiled by Claude Sonnet 5 (claude-sonnet-5) for Andr\'as Kornai}
\date{}
\begin{document}
\maketitle
\section{Introduction}
Truk\'a (ISO 639-3 tka) was the language of the Truk\'a people of the S\~ao Francisco River, in the Ilha da Assun\c{c}\~ao area near Cabrob\'o, Pernambuco, Brazil. It is unattested and unclassified. The searches found no vocabulary or word list.
\section{Primary Educational Resources}
None found.
\section{Linguistic and Reference Materials}
\begin{itemize}
\item Wikipedia, ``Truk\'a language''. \url{https://en.wikipedia.org/wiki/Truk%C3%A1_language}
\end{itemize}
\section{Religious and Community Resources}
\begin{itemize}
\item Terra Ind\'igena Truk\'a (Instituto Socioambiental). \url{https://ti.socioambiental.org/en/terras-indigenas/6151}
\item Brasil de Fato interview with a Truk\'a chief on preserving traditions. \url{https://www.brasildefato.com.br/2021/05/10/foi-a-resposta-necessaria-para-a-preservacao-de-nossas-tradicoes-diz-cacique-truka/}
\end{itemize}
\section{Digital and Audio Resources}
None found.
\section{Curriculum Development Context}
\begin{itemize}
\item UFRGS article (\emph{Espa\c{c}o Ind\'igena}) on Truk\'a teachers and practices for maintaining and revitalizing native languages in Truk\'a schools. It cites Hohenthal on Kepi\'a and Dzubuku\'a, which may not be Truk\'a itself. \url{https://seer.ufrgs.br/EspacoAmerindio/article/download/117413/84738/523545}
\end{itemize}
A heritage curriculum would have to start from the community's own cultural and language practices. No documented language material was found for it to build on.
\end{document}